\section{简介}

\begin{frame}
  \frametitle{Beamer 特点：便于排版数学公式}
  
  基本上，\LaTeX 的优点{\footnotesize （和缺点）}按原样适用。
  \begin{itemize}
    \okitem \emph{编写公式比 GUI 应用程序容易得多}
    \okitem \LaTeX 容易与文档直接关联
    \ngitem \LaTeX 作图难度大一些\footcite{tikz-nakata}，编译也会慢一点
  \end{itemize}

  \begin{block}{然而, }
    研讨会中最常见的是「\Emph{定义是什么？}」
    \begin{itemize}
      \thusitem \emph{能写数学公式真是太好了}
    \end{itemize}
  \end{block}
\end{frame}

\begin{frame}
  \frametitle{使用方法：Fork \(\rightarrow\) make}

  \begin{block}{准备}
    \begin{enumerate}
    \item
      先 \emph{Fork} 这个仓库
      \begin{itemize}
        \item \(\rightarrow\) 在右上角选择你自己作为fork目的地
      \end{itemize}
    \item
      在本地构建时，\\
      除了正常的 \LaTeX 环境 \emph{安装有 latexmk}
      \begin{itemize}
        \butitem 如果没有，请安装 docker
      \end{itemize}
    \end{enumerate}
  \end{block}

  \begin{block}{构建 pdf}
    使用 latexmk 时\ \raise1pt\hbox{\fbox{\lstinline[language = sh]|make|}}
  \begin{itemize}
    \butitem
    使用docker 时\ \raise1pt\hbox{\fbox{\lstinline[language = sh]|make ondocker|}}
  \end{itemize}
  \end{block}
\end{frame}

\section{第一张幻灯片}

\begin{frame}[fragile]
  \frametitle{第一张幻灯片：简单的示例}
  \begin{columns}
    \begin{column}{0.55\textwidth}
      \lstinputlisting[
        language = tex,
        linewidth = \textwidth
      ]{./fig/zh-cn/sample-02.tex}
    \end{column}
    \begin{column}{0.45\textwidth}
      \begin{itemize}
        \item
        整体正文字符的大小\\
        可以指定{\footnotesize （现在 10 pt）}
        \item
          注释 \alert{以 \% 开始}
        \item
          当你想腾出空间时\\
          「\textbackslash\textvisiblespace」
        \item
        Beamer 是什么，\\
        一张一张的幻灯片\\
        称为「\Emph{帧}」
      \end{itemize}
    \end{column}
  \end{columns}

\end{frame}

\begin{frame}[fragile]
  \frametitle{第一张幻灯片：项目符号/列表}
  \begin{columns}
    \begin{column}{0.51\textwidth}
      \lstinputlisting[
        language = tex,
        linewidth = \textwidth
      ]{./fig/zh-cn/sample-03.tex}
    \end{column}
    \begin{column}{0.49\textwidth}
      使用 itemize 表示项目符号
  
      \begin{itemize}
      \item 第一个项目
      \item 第二个项目
      \item 第三个项目
      \end{itemize}
    \end{column}
  \end{columns}

  \begin{columns}
    \begin{column}{0.51\textwidth}
      \lstinputlisting[
        language = tex,
        linewidth = \textwidth
      ]{./fig/zh-cn/sample-04.tex}
    \end{column}
    \begin{column}{0.49\textwidth}
      枚举型的 顺序 列表 enumerate
  
      \begin{enumerate}
      \item 第一个项
      \item 第二个项
      \end{enumerate}
    \end{column}
  \end{columns}

\end{frame}

\begin{frame}[fragile]
  \frametitle{第一张幻灯片：提高要点的{\normalsize 一点}技巧}
  \begin{columns}
    \begin{column}{0.54\textwidth}
      \lstinputlisting[
        language = tex,
        linewidth = \textwidth
      ]{./fig/zh-cn/sample-05.tex}
    \end{column}
    \begin{column}{0.46\textwidth}
      \begin{itemize}
        \ngitem
        如果这都是要点，\\
        我不明白项目之间的关系
        \thusitem 用好符号可能会很好
      \end{itemize}
    \end{column}
  \end{columns}

  \begin{mybox}{\texttt{\textbackslash item} 宏列表%
        {\footnotesize （任意定义）}
    }
    \begin{itemize}
      \okitem \verb|\okitem|: OK !
      \ngitem \verb|\ngitem|: NG
      \thusitem \verb|\thusitem|: Thus, \dots
      \butitem \verb|\butitem|: But, \dots
      \egitem \verb|\egitem|: For example, \dots
    \end{itemize}
  \end{mybox}
\end{frame}

\begin{frame}[fragile]
  \frametitle{第一张幻灯片：强调}
  \begin{enumerate}
  \item
    \emph{強調} 这样做 \verb|\emph{強調}|%
    \footnote{\insertfootnotemark
    我在未经许可的情况下覆盖了定义
    }
  \item
    \Emph{高度强调} 这样做 \verb|\Emph{高度强调}|%
    \footnote{\insertfootnotemark
    我随意定义的
    }
  \item \EMPH{疯狂强调} 可以这样做 \verb|\EMPH{疯狂强调}|%
    \footnote{\insertfootnotemark
    我随意定义的
    }
  \end{enumerate}
  
\end{frame}


\begin{frame}[fragile]
  \frametitle{第一张幻灯片:块・定义・示例}
  \begin{columns}
    \begin{column}{0.52\textwidth}
      \lstinputlisting[
        language = tex
      ]{./fig/zh-cn/sample-06.tex}
    \end{column}
    \begin{column}{0.48\textwidth}
      \begin{block}{块标题}
        最初，块填充了与背景不同的颜色,
        看起来很乱所以我停下来了
      \end{block}
    \end{column}
  \end{columns}
  \begin{columns}
    \begin{column}{0.52\textwidth}
      \lstinputlisting[
        language = tex
      ]{./fig/zh-cn/sample-07.tex}
    \end{column}
    \begin{column}{0.48\textwidth}
      \begin{definition}
        只有 1 和自身有除数，\\
        且大于 1 的自然数称为 \emph{素数}
      \end{definition}

      \begin{example}   
        \begin{itemize}
          \okitem 2 是素数．
          \okitem 3 也是素数．
          \ngitem 4 不是质数．
        \end{itemize}
      \end{example}
    \end{column}
  \end{columns}
\end{frame}

\begin{frame}[fragile]{第一张幻灯片：源代码}
  \lstinputlisting[
    language = tex,
    linewidth = 0.7\textwidth
  ]{./fig/zh-cn/sample-10.tex}
  \begin{itemize}
    \item
    frame 输入代码的时候 \Emph{需要开启 fragile 选项}
    \item
    使用相对 Makefile 所在位置的相对路径指定源文件。\\
    （在这个例子中，\raise1pt\hbox{\fbox{\lstinline[language=sh]|./fig/zh-cn/sample-10.tex|}}）
    \item
    有多种选择
    \begin{itemize}
      \egitem language, linewidth, $\dots$
    \end{itemize}
  \end{itemize}

\end{frame}

\begin{frame}[fragile]
  \frametitle{第一张幻灯片：插入图片}
  \begin{columns}
    \begin{column}{0.67\textwidth}
      \lstinputlisting[
        language = tex,
        linewidth = \textwidth
      ]{./fig/zh-cn/sample-11.tex}
    \end{column}
    \begin{column}{0.33\textwidth}
      \includegraphics[width = 3cm]{./fig/logo}
    \end{column}
  \end{columns}
    该文件由 Makefile 所在位置的相对路径指定。\\
    （这是样例，\raise1pt\hbox{\fbox{\lstinline[language=sh]|./fig/logo.jpg|}}）  
\end{frame}

\section{一点点定制}

{
\definecolor{maincolor}{RGB}{210, 100, 130}
\definecolor{accentcolor}{RGB}{0, 170, 210}

\begin{frame}{一点点定制：改变颜色主题}
  \begin{block}{颜色变更}
    \Emph{maincolor}，\EMPH{accentcolor} 设置的是，\\
    更改整个幻灯片的颜色
    \lstinputlisting[
      language = tex,
      linewidth = \textwidth
    ]{./fig/zh-cn/sample-13.tex}
    \begin{itemize}
      \ngitem \textcolor{blue}{原色} 不像现在了
      \okitem 尝试找一个冷色调
    \end{itemize}
  \end{block}

\end{frame}
}

{
\setbeamertemplate{headline}{}
\setbeamercolor{footline}{fg = backgroundcolor}

\begin{frame}{一点点定制：标题和进度表}
\begin{block}{标题}
    在幻灯片的顶部是 \emph{标题}
    \lstinputlisting[
      language = tex,
      linewidth = \textwidth
    ]{./fig/zh-cn/sample-14.tex}
    你写的时候会消失
    \begin{itemize}
    \okitem
      如果你在幻灯片开始之前写\\
      没有浪费的空间
    \end{itemize}
\end{block}

\begin{block}{进度表}
    在幻灯片的底部是 \emph{进度表}
    \lstinputlisting[
      language = tex,
      linewidth = \textwidth
    ]{./fig/zh-cn/sample-12.tex}
    如果你写，它会消失（与背景色相同）
  \end{block}

\end{frame}
}

